% TOP_PROBLEM: {"id": 3115, "title": "Star chromatic index of cubic graphs", "category_id": 3, "category": {"id": 3, "name": "graph_theory", "display_name": "Graph Theory", "description": "Problems involving graphs, networks, and their properties.", "slug": "graph-theory", "order_index": 3, "created_at": "2026-07-31T15:26:25.671Z"}, "status": "open", "problem_number": "OPG-37271", "set_id": 12}
% FIRST_POSTED: 2026-10-01
% ENTRY_KIND: statement_only
% UnsolvedMath ID 3115; source catalogue code OPG-37271
% !TeX program = lualatex
% Definition and sources only; this file is not an LLM solution attempt.
\documentclass[11pt,a4paper]{article}
\usepackage[margin=25mm]{geometry}
\usepackage{fontspec}
\setmainfont{lmroman10-regular.otf}[BoldFont=lmroman10-bold.otf,
 ItalicFont=lmroman10-italic.otf,BoldItalicFont=lmroman10-bolditalic.otf]
\usepackage{amsmath,amssymb,mathtools}
\usepackage{unicode-math}
\setmathfont{latinmodern-math.otf}
\AtBeginDocument{\let\setminus\smallsetminus}
\usepackage{xurl,hyperref}
\hypersetup{colorlinks=true,urlcolor=blue}
\setcounter{secnumdepth}{0}
\setlength{\parindent}{0pt}
\setlength{\parskip}{5pt}
\setlength{\emergencystretch}{3em}
\begin{document}
\section*{Star chromatic index of cubic graphs}\label{3115}
{\small UnsolvedMath ID 3115 · OPG-37271 · Graph Theory · OpenGarden}
\subsection{Definitions and mathematical statement}
Let $G=(V,E)$ be a finite simple undirected graph. Its maximum degree
$\Delta(G)$ is the largest number of edges incident with a vertex.
It is subcubic if $\Delta(G)\le3$, and cubic if every vertex has degree
three. An edge colouring $\varphi:E\to\{1,\ldots,q\}$ is proper when
edges with a common endpoint have different colours.

Such a colouring is a star edge colouring if no path of four edges and
no cycle of four edges uses only two colours. A path here has five
distinct vertices; it need not be an induced path. In a proper colouring,
a forbidden two-coloured four-edge path or cycle necessarily alternates
between its two colours. Equivalently, the subgraph consisting of any
two colour classes has components that are paths of at most three edges,
together with isolated vertices.

The star chromatic index $\chi'_s(G)$ is the least number of colours
in a star edge colouring. The conjecture is
\[
 \chi'_s(G)\le6\qquad\text{for every finite subcubic graph }G.
\]
The bound is independent of the number of vertices and of whether $G$
is connected, planar or bipartite. There is no preassigned colouring or
list of allowed colours at individual edges.

The restriction concerns paths with four edges, not four vertices, and
applies even when extra edges join vertices of the path. Ordinary proper
edge colouring alone does not impose it. A resolution must either give
the six-colour guarantee for every such graph or exhibit a subcubic
graph for which every proper six-colouring has a forbidden configuration.
\subsection{Short English statement}
Can every graph of maximum degree at most three be edge-coloured with six colours while avoiding every two-coloured four-edge path or cycle?
\subsection{Sources}
\begin{itemize}
\item[S1] ULAM.AI and the UnsolvedMath Contributors, supplied problems.json, record 3115 (OPG-37271), OpenGarden collection. Catalogue identity and source transcription; CC BY 4.0.
\url{https://huggingface.co/datasets/ulamai/UnsolvedMath}.
\item[S2] Open Problem Garden, Star chromatic index of cubic graphs. Original problem and discussion; the mathematical formulation below is expanded from the supplied source record.
\url{http://www.openproblemgarden.org/op/star_chromatic_index_of_cubic_graphs}.
\end{itemize}
\end{document}
